
\section{Einleitung}
Im Elektrotechnik-Labor werden grundlegende Prinzipien von Gleich- und Wechselstromkreisen untersucht. Ziel ist es, die Eigenschaften dieser Schaltungen zu analysieren, ihr Verhalten unter verschiedenen Bedingungen zu beobachten und präzise Messungen mit spezifischen Geräten durchzuführen. Der Schwerpunkt liegt dabei auf der praktischen Anwendung der theoretischen Grundlagen.

Einige der theoretischen Grundlagen sowie ergänzende Informationen zur Analyse werden aus Fachliteratur entnommen, insbesondere aus den Werken von Marinescu et al. (\textit{Elektrotechnik für Studium und Praxis} \cite{marinescu2021}, \textit{Elektrische und magnetische Felder} \cite{marinescu2012}) sowie Böhmer et al. (\textit{Elemente der angewandten Elektronik} \cite{böhmer2114}).

Im Folgenden sind die für den Versuch verwendeten Geräte aufgelistet:

\begin{itemize}
    \item \textbf{PC-Electronic-Board (hps Systemtechnik GmbH)}: Lern- und Experimentiersystem für die Elektrotechnik.
    \item \textbf{Digitalmultimeter Amprobe AM-510-EUR}: Zur präzisen Messung von elektrischen Größen wie Spannung, Strom und Widerstand.
    \item \textbf{Frequenzgenerator AIM-TTi TG1000}: Zur Erzeugung von Signalen unterschiedlicher Frequenz und Amplitude.
    \item \textbf{Keysight DSOX1102A}: Zweikanal-Oszilloskop zur Darstellung und Analyse von Signalverläufen.
    \item \textbf{Bananenstecker}: Zur sicheren und flexiblen Verbindung der Schaltungskomponenten.
\end{itemize}

Die für die Versuche verwendeten Schaltbilder und Aufbauanleitungen wurden der \textit{Praktikumsanleitung Labor Elektrotechnik} \cite{praktikumsanleitung2024} entnommen und direkt übernommen, um eine einheitliche Dokumentation zu gewährleisten.

Der \textbf{vollständige rechnerische Weg} zur Überprüfung der Ergebnisse ist im Anhang~\ref{appA} detailliert beschrieben.

\newpage
